\chapter{Thảo luận, hạn chế và kết luận}\label{chap:9}

\section{Những quan sát chính}

\begin{enumerate}
  \item \textbf{Mốc chính xác thay đổi cách đọc kết quả.} Thưởng đánh giá lấy mẫu dao động theo chuỗi nhu cầu; giá trị
    chính xác của chính sách thì không. Ở baseline, RL đạt $\SL{ch6_a/eval_mean}$ trên $10.000$ bước, đúng bằng tác tử luật
    trên cùng chuỗi nhu cầu, và giá trị chính xác của cả hai là $\SL{ch6_a/exact}$: hai con số khác $35$ chỉ do may rủi
    của mẫu.
  \item \textbf{``RL tốt hơn luật'' không phải là một mệnh đề chung.} Ở $\delta=20\%$ luật đã tối ưu, nên RL tốt nhất cũng
    chỉ hòa; RL trung bình còn thua nhẹ vì có seed học sai (chỉ $\SL{ch6_e/agg/n_opt}/10$ seed đạt tối ưu). Ở $\delta=5\%$, RL thắng.
    Ở $\delta=10\%$, RL lại thua rõ (KTC $[\SL{ch6_f/010/ci_lo};\,\SL{ch6_f/010/ci_hi}]$) vì nó giảm giá ở nhu cầu thấp,
    tham lợi trước mắt, mà không thấy rằng chính hành động đó kích hoạt sự trả đũa của đối thủ khi nhu cầu cao.
  \item \textbf{Trạng thái quan trọng hơn thuật toán.} Thêm $a_i(t-1)$ vào trạng thái giảm regret ở ba trên năm
    kịch bản của câu~10; học mô hình đối thủ rồi lập kế hoạch thì xử lý trọn lỗi trên: tác tử dựa trên mô hình có regret
    bằng $0$ ở cả năm kịch bản.
  \item \textbf{Siêu tham số quyết định độ ổn định.} $\alpha=0{,}50$ làm chính sách tham lam đổi trung bình
    $\SL{ch6_c/0.50/switches}$ lần trong huấn luyện; không khám phá ($\varepsilon=0$) chỉ đạt tối ưu ở
    $\SL{ch6_d/khong/n_opt}/10$ seed.
  \item \textbf{Phần thưởng phải khớp mục tiêu kinh doanh.} Huấn luyện theo doanh thu làm lợi nhuận thật tụt còn
    $\SL{ch6_g/doanh_thu/exact_mean}$ mỗi bước.
\end{enumerate}

\section{Hạn chế}

Mô hình cố ý đơn giản: hai mức nhu cầu, hai hành động, không tồn kho, không mùa vụ, khách hàng không nhớ giá. Kết quả chỉ
có giá trị trong phạm vi mô hình và không nên suy rộng thành kết luận kinh tế. Mười seed là đủ để thấy độ phân tán nhưng
khoảng tin cậy vẫn rộng. Đề gốc và mã tham chiếu không có sẵn trên máy làm bài, nên mô hình được dựng lại từ công thức
thống nhất của đề; mọi thay đổi về cách đo được ghi ở Chương~\ref{chap:5}. Hệ mua sắm ở Chương~\ref{chap:8} dùng dữ liệu truy vấn
tổng hợp và lời gọi đồng bộ trong một tiến trình.

\section{Kết luận}

Báo cáo đã mô hình hóa, cài đặt và đánh giá hai cách ra quyết định trong một hệ đa tác tử thương mại điện tử. Ba điểm rút
ra: đánh giá phải tách khỏi huấn luyện và nên có mốc tối ưu chính xác; ưu thế của học hay của luật phụ thuộc vào thiết kế
phần thưởng và vào việc trạng thái có đủ thông tin Markov hay không; và một tác tử dựa trên mô hình nhỏ, biết đối thủ
phản ứng ra sao, vượt cả Q-learning lẫn luật cố định. Hướng tiếp theo là thêm tồn kho và nhiều mức giá, đối thủ có trí nhớ
dài hơn, và đánh giá các thuật toán đa tác tử như học đồng thời có mô hình đối thủ.
